\section{Data and manipulation-simulation plan}\label{sec:data}
\begin{keybox}
\textbf{In one sentence.} Authentic public Sentinel-2 patches (EuroSAT) are split without leakage; every manipulation is generated by deterministic, seeded code, so for every patch we know whether it was altered, how, and where, and the methods used for testing are kept apart from those used for training.
\end{keybox}

\subsection{Data sources}
\begin{table}[!htbp]
\centering\small
\begin{tabularx}{\textwidth}{|L{3.6cm}|Y|L{3.6cm}|}
\hline
\textbf{Source} & \textbf{What we take} & \textbf{Used for}\\
\hline
EuroSAT RGB \cite{eurosat} & 27,000 Sentinel-2 patches of 64$\times$64 pixels, 10 land-cover classes, MIT licence; fetched by \cd{scripts/download_eurosat.py}, which refuses data whose class counts differ & authentic patches; every split and every number in this report\\
\hline
\cd{samples/authentic/} & 30 unmodified official patches shipped in the repository & labs, tests and the demo run without the 2~GB dataset\\
\hline
Emulated source and processor & signing keys \cd{sat-01} and \cd{proc-01}, capture claims, signed processing steps (software keystore) & provenance of the 11-attack campaign and the labs\\
\hline
Jury patches & held-out authentic and manipulated patches supplied by the organisers (any format or size) & the standard scenario; read by \cd{pixelproof.cli detect}\\
\hline
\end{tabularx}
\caption{Data sources. No private or restricted data is used; no paid service is needed.}
\end{table}

\subsection{Splits and leakage rules}
Every official EuroSAT patch is assigned to exactly one split by a SHA-256 hash of its id: \emph{train} 60\%, \emph{validation} 15\%, \emph{test-seen} 12.5\% and \emph{test-held-out} 12.5\%. For each source patch we render one authentic and one manipulated copy. A patch, its manipulated copy and its splicing donor always share a split, so there is \textbf{no leakage}. Thresholds and the fusion weight are fixed on validation only. The seen test set holds {{m:/models/fused (w=0.4)/test_dev/n|int}} patches and the held-out test set {{m:/models/fused (w=0.4)/test_heldout/n|int}}. The full data card is \cd{docs/DATASET_CARD.md}.

\begin{figure}[!htbp]
\centering
\begin{adjustbox}{max width=\textwidth}
\begin{tikzpicture}[font=\scriptsize]
  \node[bg,text width=27mm] (a) at (0,0) {\textbf{1. EuroSAT RGB}\\27,000 official patches};
  \node[bg,text width=27mm] (b) at (3.5,0) {\textbf{2. Split}\\SHA-256 of the patch id};
  \node[bt,text width=27mm] (c) at (7,0) {\textbf{3. Generate}\\one authentic + one manipulated copy, seeded};
  \node[bo,text width=27mm] (d) at (10.5,0) {\textbf{4. Fit and calibrate}\\train 60\%, validation 15\%: thresholds for the budget};
  \node[bn,text width=27mm] (e) at (14,0) {\textbf{5. Score}\\seen test, held-out v1, frozen v2};
  \draw[arr] (a)--(b); \draw[arr] (b)--(c); \draw[arr] (c)--(d); \draw[arr] (d)--(e);
  \node[lbl,align=center,text width=40mm] at (7,-1.3) {every evaluation patch is rebuilt deterministically from (patch id, split, family)};
\end{tikzpicture}
\end{adjustbox}
\caption{Data pipeline: from the official release to scored results; each step is a script in \texttt{scripts/}.}
\end{figure}

\subsection{Manipulation generators: seen, held-out and frozen}
The standard scenario asks us to detect a resampled or spliced image, a copy-move alteration and a generated or inpainted region. Our generators cover these three groups with methods that are \textbf{different on purpose} between training and testing, so the held-out numbers measure generalisation, not memorisation.

\begin{table}[!htbp]
\centering\small
\begin{tabularx}{\textwidth}{|L{2.5cm}|Y|Y|Y|}
\hline
\textbf{Group} & \textbf{Training (seen)} & \textbf{Held-out v1} & \textbf{Held-out v2 (frozen before any change)}\\
\hline
Resampled / spliced & bilinear $\times$1.15--1.5, rotation 4--12$^\circ$, hard-edged splices & bicubic/Lanczos, $\times$1.6--2.2, rotation 15--35$^\circ$, feathered colour-matched splices & perspective/biquadratic warp; Poisson (seamless-clone) splices, donor from another land-cover class\\
\hline
Copy-move & translated rectangles & rotated/rescaled/noisy polygons & mirrored elliptical region, soft-blended\\
\hline
Generated / inpainted & Telea inpainting & Navier--Stokes, biharmonic, texture synthesis & exemplar fill + injected sensor noise\\
\hline
Laundering (all patches) & 1--2 of JPEG/WebP/noise & JPEG 2000 + radiometric stretch & ---\\
\hline
\end{tabularx}
\caption{Manipulation methods per evaluation set.}\label{tab:methods}
\end{table}

\textbf{Freezing protocol.} Held-out v2 was generated and its SHA-256 recorded before any detector change made after v1:
\begin{center}\ttfamily\footnotesize {{h:/heldout_v2_sha256|raw}}\end{center}
The v1-model baseline on it is stored next to the result of the shipped model, so the gain from our changes is honest. A third set, held-out v3 (thin-strip regions; Fourier anisotropic resampling, Laplacian-pyramid strip splicing \cite{burt}, multi-clone strip copy-move, push-pull multiscale fill with grain), was generated from the seen-test split with methods that appear nowhere else and frozen before the Phase 2 candidate was trained. It holds {{v3:/n_patches|int}} patches ({{v3:/counts/authentic|int}} authentic, {{v3:/counts/resample|int}} resampled, {{v3:/counts/splice|int}} spliced, {{v3:/counts/copy_move|int}} copy-move, {{v3:/counts/inpaint|int}} inpainted); its SHA-256 is
\begin{center}\ttfamily\footnotesize {{v3:/sha256|raw}}\end{center}
The scoring script refuses to score the same model on v3 twice.

\subsection{What is real and what is emulated}
\begin{itemize}
  \item \textbf{Real:} the pixels (official Sentinel-2 patches), the cryptography (Ed25519, SHA-256 from a standard library), the network attacks (real packets through a real relay or router), the detector.
  \item \textbf{Emulated:} the satellite and its processing chain (a Python process holding two software keys), the capture claim fields, and the manipulations of unsigned imagery (generated by our own code, not by third-party generative models).
\end{itemize}
